% XeLaTeX can use any Mac OS X font. See the setromanfont command below.
% Input to XeLaTeX is full Unicode, so Unicode characters can be typed directly into the source.

% The next lines tell TeXShop to typeset with xelatex, and to open and save the source with Unicode encoding.

%!TEX TS-program = xelatex
%!TEX encoding = UTF-8 Unicode

\documentclass[11pt,twoside,openany]{book}
\usepackage{multirow}
\usepackage{pdfpages}

\usepackage{xspace}

\usepackage{array}
\usepackage{lipsum}

\usepackage{algorithm}
\usepackage{algpseudocode}

\usepackage{geometry}                % See geometry.pdf to learn the layout options. There are lots.
%\usepackage[margin=1cm]{geometry}


% \addtolength{\textwidth}{1.75in}
%
% \addtolength{\topmargin}{-.875in}
% \addtolength{\textheight}{1.75in}

\geometry{a4paper,left=20mm,top=20mm,bottom=50mm,total={162mm,230mm}}                   % ... or a4paper or a5paper or ...
%\geometry{landscape}                % Activate for for rotated page geometry
%\usepackage[parfill]{parskip}    % Activate to begin paragraphs with an empty line rather than an indent
\usepackage{amssymb}
%\usepackage{todonotes}
\setlength{\marginparwidth}{2cm}
\usepackage[backgroundcolor=white,bordercolor=blue,linecolor=blue,textwidth=1cm]{todonotes}
\usepackage{booktabs}
\usepackage{longtable}
\usepackage{url}

\usepackage[depth=4]{bookmark}


\usepackage{algorithm}
\usepackage{algpseudocode}

\usepackage{imakeidx}
\usepackage[utf8]{inputenc}
\usepackage[T1]{fontenc}
\makeindex[intoc]

\usepackage[font=normalsize]{caption}

\renewcommand{\theparagraph}{\S\arabic{paragraph}}
\setcounter{secnumdepth}{4}


\usepackage{xcolor}
\usepackage{listings}
\usepackage{amsmath}

\makeatletter
\global\let\tikz@ensure@dollar@catcode=\relax
\makeatother

\definecolor{mGreen}{rgb}{0,0.6,0}
\definecolor{mGray}{rgb}{0.5,0.5,0.5}
\definecolor{mPurple}{rgb}{0.58,0,0.82}
\definecolor{backgroundColour}{rgb}{0.95,0.95,0.92}
\definecolor{gray}{rgb}{0.4,0.4,0.4}
\definecolor{darkblue}{rgb}{0.0,0.0,0.6}
\definecolor{cyan}{rgb}{0.0,0.6,0.6}

\lstdefinestyle{CStyle}{
    backgroundcolor=\color{backgroundColour},
    commentstyle=\color{mGreen},
    keywordstyle=\color{magenta},
    numberstyle=\tiny\color{mGray},
    stringstyle=\color{mPurple},
    basicstyle=\scriptsize\ttfamily,
    breakatwhitespace=false,
    breaklines=true,
    captionpos=b,
    keepspaces=true,
    numbers=left,
    numbersep=5pt,
    showspaces=false,
    showstringspaces=false,
    showtabs=false,
    tabsize=2,
    language=C
}


\lstset{
  backgroundcolor=\color{white},   % choose the background color; you must add \usepackage{color} or \usepackage{xcolor}; should come as last argument
  basicstyle=\footnotesize\ttfamily,        % the size of the fonts that are used for the code
  breakatwhitespace=false,         % sets if automatic breaks should only happen at whitespace
  breaklines=true,                 % sets automatic line breaking
  captionpos=b,                    % sets the caption-position to bottom
  commentstyle=\color{gray},    % comment style
  deletekeywords={...},            % if you want to delete keywords from the given language
  %escapeinside={\%*}{*)},          % if you want to add LaTeX within your code
  %extendedchars=true,              % lets you use non-ASCII characters; for 8-bits encodings only, does not work with UTF-8
  %firstnumber=1000,                % start line enumeration with line 1000
  frame=single,                    % adds a frame around the code
  keepspaces=true,                 % keeps spaces in text, useful for keeping indentation of code (possibly needs columns=flexible)
  keywordstyle=\color{blue},       % keyword style
  language=C,                 % the language of the code
  %morekeywords={*,...},            % if you want to add more keywords to the set
  numbers=left,                    % where to put the line-numbers; possible values are (none, left, right)
  numbersep=5pt,                   % how far the line-numbers are from the code
  numberstyle=\tiny\color{gray}, % the style that is used for the line-numbers
  rulecolor=\color{black},         % if not set, the frame-color may be changed on line-breaks within not-black text (e.g. comments (green here))
  showspaces=false,                % show spaces everywhere adding particular underscores; it overrides 'showstringspaces'
  showstringspaces=false,          % underline spaces within strings only
  showtabs=false,                  % show tabs within strings adding particular underscores
  stepnumber=1,                    % the step between two line-numbers. If it's 1, each line will be numbered
  stringstyle=\color{black},     % string literal style
  tabsize=2,                     % sets default tabsize to 2 spaces
}

\lstset{
  columns=fullflexible,
  showstringspaces=false,
  commentstyle=\color{gray}\upshape,
  backgroundcolor=\color{backgroundColour},
  commentstyle=\color{mGreen},
  keywordstyle=\color{magenta},
  numberstyle=\tiny\color{mGray},
  stringstyle=\color{mPurple},
  basicstyle=\scriptsize\ttfamily,
  tabsize=2
}

\lstdefinelanguage{XML}
{
  morestring=[b]",
  morestring=[s]{>}{<},
  morecomment=[s]{<?}{?>},
  stringstyle=\color{black},
  identifierstyle=\color{darkblue},
  keywordstyle=\color{cyan},
  morekeywords={xmlns,version,type}% list your attributes here
}


% Will Robertson's fontspec.sty can be used to simplify font choices.
% To experiment, open /Applications/Font Book to examine the fonts provided on Mac OS X,
% and change "Hoefler Text" to any of these choices.

\usepackage{fontspec,xltxtra,xunicode}
\defaultfontfeatures{Mapping=tex-text}
\setromanfont[Mapping=tex-text]{Optima}
\setsansfont[Scale=MatchLowercase,Mapping=tex-text]{Optima}
\setmonofont[Scale=MatchLowercase]{Andale Mono}

%%%% Graphics
\usepackage{graphicx}
\usepackage{adjustbox}
\usepackage{tikz}
\definecolor{unilublue}{RGB}{55,149,218}
\definecolor{sntred}{RGB}{219,46,27}
\definecolor{sntpurple}{RGB}{86,30,130}
%\definecolor{sntblue}{RGB}{50,130,207}
\definecolor{sntblue}{RGB}{55,149,218}

\pgfdeclareimage[width=30mm]{logo-snt}{logos/logo-snt}
\pgfdeclareimage[width=30mm]{logo-uni-lu}{logos/logo-uni-lu}

%%%% Fancy
\usepackage{fancyhdr}

%%%% Increase page length
\addtolength{\textheight}{1in}

%%%% Sections
\usepackage{sectsty}
\allsectionsfont{\sffamily}

\setlength{\parskip}{1em}

%\renewcommand{\,}{$^{\cdot}$}

\usepackage{hyperref}
\hypersetup{
    colorlinks,
    citecolor=black,
    filecolor=black,
    linkcolor=black,
    urlcolor=black
}


%%%% Title
\newcommand{\titleone}{\textsf{FAQAS Framework - DAMAt}}
\newcommand{\titletwo}{\textsf{Verification Report}}
\newcommand{\titlethree}{\textsf{ESAIL-ADCS Case Study}}

\newcommand{\todoinline}[1]{\todo[color=orange,inline]{ \textbf{TODO}: #1 }}

\newcommand{\TODO}[1]{\todo[color=orange,inline]{ \textbf{TODO}: #1 }}
\newcommand{\OSCAR}[1]{\todo[color=yellow,inline]{ \textbf{OSCAR}: #1 }}
\newcommand{\DONE}[1]{\todo[color=green,inline]{ \textbf{DONE}: #1 }}

\newcommand{\CHANGEDTWO}[1]{\textcolor{black}{#1}}

\newcommand{\EMPH}[1]{\textbf{\emph{#1}}}
\newcommand{\INDEX}[1]{\index{\MakeLowercase{#1}}\EMPH{#1}}

\newcommand{\MASS}{\textit{MASS}\xspace}

\newcommand{\DAMA}{\textit{DAMAt}\xspace}

\newcommand{\SEMUS}{\textit{SEMuS}\xspace}

\newcommand{\SEMU}{\textit{SEMu}\xspace}

\newcommand{\case}{\textit{ESAIL-ADCS case study}\xspace}

\newcommand{\D}{$\Delta$\xspace}

\newcommand{\remark}{\textbf{Remark:}\xspace}

\newcounter{paranum}
\newcommand{\action}{\vspace{10pt}\noindent\textbf{Action Item \refstepcounter{paranum}\theparanum: }}

\begin{document}

\pagestyle{fancy}
\renewcommand{\sectionmark}[1]{\markright{\textit{#1}}}

\renewcommand{\headrulewidth}{2pt}% 2pt header rule
\renewcommand{\headrule}{\hbox to\headwidth{%
  \color{sntblue}\leaders\hrule height \headrulewidth\hfill}}

\fancyhf{}

%\lhead{\fancyplain{}{\setlength{\unitlength}{1mm}
%\begin{picture}(0,0)
%\put(0,-4){\includegraphics[width=50pt]{logos/logo-snt}}
%\end{picture}}}

\newcommand{\DOCUMENTID}{Verification Report: ESAIL-ADCS}
\lhead{\fancyplain{}{\textit{\DOCUMENTID}}}
\rhead{\fancyplain{}{\rightmark }}


\fancyfoot[C]{%
\begin{tikzpicture}[remember picture,overlay]
\path [fill=sntred]    ([xshift=88pt,yshift=20pt]current page.south west) rectangle
                       ([xshift=229pt,yshift=50pt] current page.south west);
\path [fill=sntpurple] ([xshift=229pt,yshift=20pt] current page.south west) rectangle
                       ([xshift=370pt,yshift=50pt] current page.south west);
\path [fill=sntblue]   ([xshift=370pt,yshift=20pt] current page.south west) rectangle
                       ([xshift=510pt,yshift=50pt] current page.south west);
\end{tikzpicture}
}

\fancyfoot[RO]{
\begin{tikzpicture}[remember picture,overlay]
\node [circle, ultra thick, fill=white, draw=sntblue] at ([xshift=530pt,yshift=35pt] current page.south west) {\thepage};
\end{tikzpicture}}

\fancyfoot[LE]{
\begin{tikzpicture}[remember picture,overlay]
\node [circle, ultra thick, fill=white, draw=sntblue] at ([xshift=65pt,yshift=35pt] current page.south west) {\thepage};
\end{tikzpicture}}


\newcommand{\STARTCHANGEDWPT}{}
\newcommand{\ENDCHANGEDWPT}{}

\newcommand{\CHANGED}[1]{{#1}}
\newcommand{\prog}{P}
\newcommand{\prop}{}

\newcommand{\FAQAS}{FAQAS-framework\xspace}


\newcommand{\REVISION}[2]{{#2}}
\newcommand{\REVTOOL}[2]{#2}

\newcommand{\REVTWO}[2]{#2}

\thispagestyle{empty}

\begin{tikzpicture}[remember picture,overlay]
\path [fill=sntred]    ([xshift=30pt,yshift=20pt]current page.south west) rectangle
                       ([xshift=210pt,yshift=50pt] current page.south west);
\path [fill=sntpurple] ([xshift=210pt,yshift=20pt] current page.south west) rectangle
                       ([xshift=390pt,yshift=50pt] current page.south west);
\path [fill=sntblue]   ([xshift=390pt,yshift=20pt] current page.south west) rectangle
                       ([xshift=570pt,yshift=51pt] current page.south west);
\path [fill=unilublue] ([xshift=30pt,yshift= 50pt] current page.south west) --
                       ([xshift=570pt,yshift= 50pt] current page.south west)
                       [rounded corners=20pt] --
                       ([xshift=570pt,yshift=740pt] current page.south west)
                       [sharp corners] --
                       ([xshift=30pt,yshift=740pt] current page.south west);
\node [fill=white,rounded corners=0pt,inner xsep=6pt,inner ysep=3pt]
      at ([xshift=523pt,yshift=120pt] current page.south west)
      {\pgfuseimage{logo-uni-lu}};

\node [fill=white,rounded corners=2pt,inner xsep=6pt,inner ysep=3pt]
      at ([xshift=520pt,yshift=780pt] current page.south west)
      {\pgfuseimage{logo-snt}};

%\node [circle, fill=white, draw=sntblue] at ([xshift=550pt,yshift=35pt] current page.south west) {a};

\node[draw=none,fill=none,right] at (-1, -7){\color{white}\LARGE\bf\titleone};
\node[draw=none,fill=none,right] at (-1, -8){\color{white}\LARGE\bf\titletwo};
\node[draw=none,fill=none,right] at (-1, -9){\color{white}\LARGE\bf\titlethree};

\node[draw=none,fill=none,right] at (-1, -12){\color{white}\Large\textsf{O. Cornejo, F. Pastore, E. Viganò}};
\node[draw=none,fill=none,right] at (-1, -13){\color{white}\Large\textsf{Interdisciplinary Centre for Security, Reliability and Trust}};
\node[draw=none,fill=none,right] at (-1, -14){\color{white}\Large\textsf{University of Luxembourg}};
\node[draw=none,fill=none,right] at (11, -16){\color{white}\textsf{ITT-1-9873-ESA-FAQAS-SUTP}};
\node[draw=none,fill=none,right] at (11, -17){\color{white}\Large\textsf{Issue 4, Rev. 1}};
\node[draw=none,fill=none,right] at (11, -18){\color{white}\Large\textsf{\today}};
\node[draw=none,fill=none,right] at (-1, -24){\color{white}\tiny\textsf{EUROPEAN SPACE AGENCY. CONTRACT REPORT.}};
\node[draw=none,fill=none,right] at (-1, -24.3){\color{white}\tiny\textsf{The work described in this report was done under ESA contract. Responsibility for the contents resides in the author or organisation that prepared it.}};
\node[draw=none,fill=none,right] at (-1, -24.8){\color{white}\tiny\textsf{The copyright in this document is vested in the University of Luxembourg.}};
\node[draw=none,fill=none,right] at (-1, -25.1){\color{white}\tiny\textsf{This document may only be reproduced in whole or in part, or stored in a retrieval system,or transmitted in any form, or by any means electronic,}};
\node[draw=none,fill=none,right] at (-1, -25.4){\color{white}\tiny\textsf{mechanical, photocopying or otherwise, either with the prior permission of the University of Luxembourg or in accordance with the terms of ESTEC Contract No. 4000128969/19/NL/AS.}};

\node[draw=none,fill=none,right] at (-1, -24){\color{white}\tiny\textsf{EUROPEAN SPACE AGENCY. CONTRACT REPORT.}}; 
\node[draw=none,fill=none,right] at (-1, -24.3){\color{white}\tiny\textsf{The work described in this report was done under ESA contract. Responsibility for the contents resides in the author or organisation that prepared it.}};
\node[draw=none,fill=none,right] at (-1, -24.8){\color{white}\tiny\textsf{The copyright in this document is vested in the University of Luxembourg.}};
\node[draw=none,fill=none,right] at (-1, -25.1){\color{white}\tiny\textsf{This document may only be reproduced in whole or in part, or stored in a retrieval system,or transmitted in any form, or by any means electronic,}}; 
\node[draw=none,fill=none,right] at (-1, -25.4){\color{white}\tiny\textsf{mechanical, photocopying or otherwise, either with the prior permission of the University of Luxembourg or in accordance with the terms of ESTEC Contract No. 4000128969/19/NL/AS.}};


\end{tikzpicture}

\newpage


% !TEX root = MAIN.tex

\section*{Revisions}
\label{sec:revisions}


\setlength\LTleft{0pt}
\setlength\LTright{0pt}
\tiny 
%@{\extracolsep{\fill}}
\begin{longtable}{|p{2cm}|p{1cm}|p{1.5cm}|p{9cm}|@{}}
\label{table:codeoperators} \\
\hline
\textbf{Issue Number}&\textbf{Date}&\textbf{Authors}&\textbf{Description}\\
\hline
ITT-1-9873-ESA-FAQAS-FR
Issue 1 Rev. 1&
October 29th, 2021&
Fabrizio Pastore, Oscar Cornejo, Enrico Viganò&
\begin{minipage}{8cm}
Initial release.
\end{minipage}
\\
\hline
ITT-1-9873-ESA-FAQAS-FR
Issue 1 Rev. 2&
November 11th, 2021&
Fabrizio Pastore, Oscar Cornejo, Enrico Viganò&
\begin{minipage}{8cm}
Added Chapter~\ref{chap:deliverables_summary} (deliverables summary).\\
Added Chapter~\ref{ch:toolset} (description of toolset package).\\
Added information about input and outputs of the tools in Chapter~\ref{chapter:methodology} (see blue text).
\end{minipage}
\\
\hline
                                                    
\end{longtable}
\normalsize

\clearpage
\clearpage

\tableofcontents

% !TEX root = MAIN.tex

\chapter{Executive Summary}

This document is a report on the issues identified by applying Data-driven Mutation Analysis on the ESAIL-ADCS case study using the \DAMA tool, which has been developed in the context of the \FAQAS.

\DAMA is a data-driven mutation analysis tool. It generates mutants by modiying the data contained in the SUT buffers throught the insertion in the source code of Mutation Probes. The mutation probes generate mutants following mutation operators defined in a Fault Model.

The code we analyzed was the ESAIL SVF test suite, which evaluates the ESAIL OBSW and in particular the test cases regarding the interaction between the ADCS (Attitude Determination and Control System) component and the OBSW (On Board Software).


\section{Terms, definitions and abbreviated terms}

\begin{itemize}
\item{FAQAS}: activity ITT-1-9873-ESA
\item{FAQAS-framework}: software system to be released at the end of WP4 of FAQAS.
\item{SUT}: Software under test, i.e, the software that should be mutated by means of mutation testing.
\item{FMC}: Fault model coverage, the percentage of fault models covered by the test suite.
\item{MOC}: Mutation operation coverage, the percentage of data items that have been mutated at least once, considering only those that belong to the data buffers covered by the test suite.
\item{MS}: Mutation Score, the percentage of mutants killed by the test suite (i.e., leading to at least one test case failure) among the mutants that target a fault model and for which at least one mutation operation was successfully performed.
\item{SVF}: Software Validation Facility.
\item{ADCS}: Attitude Control and Determination Software.

\end{itemize}

\clearpage


% !TEX root = MAIN.tex

\chapter{Configuration of the FAQAS Toolset}

% \TODO{that we used}



We configured MASS and SEMuS by following the default configuration provided in the Software User Manual (SUM) document, except for a couple of additional configurations, which are described in this chapter.

\section{MASS Configurations}


Concerning the configuration of the \textbf{MASS tool}, the additional configurations that we considered were the ones regarding the disregarding of equivalent and duplicate mutants with trivial compiler optimizations, and the one concerning the use of sampling.
In the case of trivial compiler optimisations we enabled all levels of optimizations permitted by GCC~\cite{GCCopt}, that is, levels \texttt{O0}, \texttt{O1}, \texttt{O2}, \texttt{O3}, \texttt{Ofast}, and \texttt{Os}. 
In the case of sampling, we decided to analyze all the mutants generated by the tool, and do not employ sampling methods. 

\section{SEMuS Configurations}

Now we report the configurations concerning the \textbf{SEMuS tool}, the additional configurations concerned the compilation command used, the timeouts, and the memory limits imposed to the test generation process.

\begin{lstlisting}[language=bash, label=libutil_com_cmd, caption=Compilation command for a single Libutil source used in SEMuS.]
-g -std=gnu99 -m64 -O0 -Wall -Wextra -Wshadow -Wcast-align -Wwrite-strings -Wno-unused-parameter -I$repo_root_dir/src -I$repo_root_dir/src/gosh -I$repo_root_dir/src/zip -I$repo_root_dir/src/watchdog  -I$repo_root_dir/build/src/watchdog -I$repo_root_dir/build/src -I$repo_root_dir/include -I$repo_root_dir/build/include -I$repo_root_dir/include/gs -I$repo_root_dir/build/include/gs -I$repo_root_dir/include/deprecated -I$repo_root_dir/build/include/deprecated -I$repo_root_dir/include/deprecated/gs/gosh -I$repo_root_dir/build/include/deprecated/gs/gosh -I$repo_root_dir/include -o $out_file $in_file
\end{lstlisting}

Listing~\ref{libutil_com_cmd} presents the compilation command we used to enable test generation on Libutil, in particular, the variable \texttt{\$repo\_root\_dir} represents the installation path of Libutil, \texttt{\$out\_file} represents the file storing the compiled object, and \texttt{\$in\_file} represents the source being analyzed.


We configured SEMuS to terminate the test generation after 300 seconds. The timeout implies that after 300 seconds, SEMuS will try to generate test inputs for the states reached so far, even if the symbolic execution has not reached the end of the function under test. If tests are generated, it means that the processed mutants have been killed. Otherwise, it means that SEMuS could not find test inputs that kill those mutants (they require manual inspection because they might be either equivalent or non-equivalent mutants).

%\DONE{What are the implication of the above? You should write that it still checks if the generated inputs kill the mutants, if not they are discarded.}

We also set SEMuS to use a maximum of 2\,000 MB for the test generation, which means that if the test generation reaches 2\,000 MB of memory consumption, the test generation process will be stopped. The limit of 2\,000 MB was selected by following the guidelines of the underlying test generation tool (i.e., SEMu~\cite{titcheu2021killing}).

%\DONE{What is the motivation for the choice above? Why do we believe is a good compromise between scalability and effectiveness?}



% !TEX root = MAIN.tex

\chapter{Results}

\section{Metrics}

Starting from the Fault Model represented in Table~\ref{tab:fault_model}, \DAMA generated 666 mutants. The mutated version of the program were executed against the \case test suite.

The results were expressed with the following three metrics:
\begin{enumerate}
\item Fault model coverage, the percentage of fault models covered by the test suite.
\item Mutation operation coverage, the percentage of data items that have been mutated at least once, considering only those that belong to the data buffers covered by the test suite.
\item Mutation Score, the percentage of mutants killed by the test suite (i.e., leading to at least one test case failure) among the mutants that target a fault model and for which at least one mutation operation was successfully performed.
\end{enumerate}


A low score in one of the metrics indicate one of following scenarios, repsectively:
\begin{enumerate}
\item The message type targeted by a fault model is never exercised.
\item The message type is covered by the test suite, but it is not possible to perform some of the mutation operations. It depends on the fact that not all the input partitions are exercised by the test suite.
\item The mutation is performed but the test suite does not fail. It may depend on two reasons: (1) the test oracles are imprecise (e.g., they do not verify all the state variables), (2) the system is not brought into a state where the effect of the mutation is notices (i.e., the scenarios exercised are insufficient).
\end{enumerate}

The results for the \case are reported in Table~

\section{Uncovered Fault Models}

\section{Uncovered Mutants}

\section{Live Mutants}


\clearpage

\printindex

\end{document}
